% VeriMem --- IEEE Transactions on Artificial Intelligence format.
%
% TAI's own class file (ieeetai.cls) is distributed only through the IEEE Author
% Center and is not on CTAN, so this uses IEEEtran with the [journal] option --
% the same two-column Transactions layout the TAI template is built on. The
% TAI-specific Impact Statement is set manually below to match the template.
%
% Build:  tectonic -X compile report/verimem_tai.tex
\documentclass[journal]{IEEEtran}

\usepackage{cite}
\usepackage{graphicx}
\usepackage{amsmath}
\usepackage{booktabs}
\usepackage{url}
\usepackage[T1]{fontenc}

% The Impact Statement block required by IEEE TAI, set to match the template.
\newenvironment{impactstatement}
  {\par\vspace{1.5ex}\noindent\textit{Impact Statement}---\ignorespaces}
  {\par\vspace{1.5ex}}

\begin{document}

\title{VeriMem: Poisoning the Memory of a\\ Retrieval-Augmented Fact-Checking Agent}

\author{
  Teesha~Ramchandani~(12342260),
  Rahul~Raj~(12341680),
  Rajeev~Kumar~(12341700),\\
  Durgesh~Nama~(12340720),
  and~Yogesh~(12342420)%
  \thanks{All authors are with the Indian Institute of Technology Bhilai,
  Chhattisgarh, India. This work was carried out as a course project in
  Natural Language Processing.}%
}

\markboth{IEEE Transactions on Artificial Intelligence, Vol.~00, No.~0, November~2026}%
{Ramchandani \MakeLowercase{\textit{et al.}}: Poisoning the Memory of a Retrieval-Augmented Fact-Checking Agent}

\maketitle

\begin{abstract}
Retrieval-augmented agents increasingly store their own completed work and
retrieve it as context for later inputs. That store is an attack surface, and
this article measures what an adversary gains from writing to it. We build
VeriMem, a fact-checking agent that reproduces two published systems from
scratch and transfers both to claim verification, then attack it at two
surfaces: the agent's memory of its own verifications, and the evidence corpus
it searches. Injecting fabricated records raises attack success from 0.220 to
0.430 as the poisoned share of memory reaches twenty percent, while an agent
with no memory stays at 0.160. Two defences, evaluator gating and Bayesian
source trust, each flatten that curve, but which one does the work inverts with
the adversary: against fabrications citing unknown publishers, trust is decisive
and gating adds nothing, whereas against an adversary citing the publishers the
agent already trusts, trust fails completely and gating carries the defence
alone. Measuring what the trust model learns explains this. It does not separate
known misinformation sources from the rest, and ranks a social network above a
scientific journal, because it records which publishers coincide with correct
verdicts rather than which publishers are reliable. It detects novelty, not
credibility. Poisoning the evidence corpus proves both more effective and less
defensible, because fabricated pages are genuinely relevant and the relevance
model scores them almost as highly as real evidence.
\end{abstract}

\begin{impactstatement}
Agents that accumulate a memory of their own work are being deployed in
settings where their judgements carry weight, including automated fact
checking. This article shows that such memory is an exploitable surface: an
attacker who can write to it steers verdicts without touching the model, the
retriever, or the evidence corpus, using records that are perfectly well formed
and so invisible to schema validation. The two defences we examine are
partial in a way practitioners should know about. Source-reputation scoring, a
common mitigation, turns out to detect unfamiliar publishers rather than
unreliable ones, which means it stops an opportunistic attacker and offers
nothing against one who impersonates trusted sources. Relevance filtering fails
in the opposite case. Deployments relying on either alone are protected against
only half the threat, and our measurements indicate which half.
\end{impactstatement}

\begin{IEEEkeywords}
Adversarial machine learning, data poisoning, fact verification,
large language models, retrieval-augmented generation, trust modelling.
\end{IEEEkeywords}

\section{Introduction}
\IEEEPARstart{A}{retrieval-augmented} agent that keeps a memory of its own past
work has a new attack surface: the memory. An adversary who can write to it
shapes how future inputs are judged without touching the model, the retriever,
or the evidence corpus. The records they write need not be malformed. A
poisoned trajectory can be a syntactically perfect record of a verification that
never happened, which is precisely what makes it hard to find: the attack is
semantic, not structural, so no schema check will catch it.

We study this concretely. VeriMem is a fact-checking agent. Given a claim it
returns a verdict in \{Supported, Refuted, Conflicting evidence, Not enough
evidence\} together with an evidence trail: the sources it used, their relevance
grades, the trust it assigns each publisher, and the past verifications it drew
on. It combines two systems we reproduce from scratch, CRAG's retrieval
evaluator and three-way routing \cite{yan2024crag}, and ExpRAG's trajectory
store feeding retrieval-augmented fine-tuning \cite{ferraz2026exprag}, and
evaluates on AVeriTeC \cite{schlichtkrull2023averitec}.

Our contributions are as follows.

\begin{enumerate}
\item A \textbf{memory-poisoning study} (Section~\ref{sec:c1}) at 0/5/10/20\%
injection, with two adversaries and two attack surfaces. Experience retrieval is
a real vulnerability, defences do cut attack success, and, the finding we
believe matters most, the two defences are \emph{complementary rather than
redundant}: each fails against the adversary the other stops.
\item \textbf{Evaluator-gated experience retrieval}, which scores a retrieved
trajectory on whether its own evidence supports its own verdict.
\item \textbf{Bayesian source trust}, together with a measurement showing it
learns something other than what it was designed to learn
(Section~\ref{sec:trust}).
\item The \textbf{first reproduction of ExpRAG} we are aware of, which does not
replicate its central claim, with the reason measured
(Section~\ref{sec:exprag}).
\end{enumerate}

Every number below carries a bootstrap confidence interval over the evaluation
set. This is not decoration. Macro-F1 on this benchmark gives a quarter of its
weight to a 38-instance class on which every system we tested scores near zero,
and our own early readings of the main table were wrong because of it.

\section{Related Work}

\subsection{Corrective Retrieval}
CRAG \cite{yan2024crag} adds a fine-tuned evaluator that grades retrieved
documents and routes on the result: refine, rewrite-and-search, or both.
Refinement is decompose-then-recompose over sentence strips. CRAG reports
consistent gains across four question-answering benchmarks.

\subsection{Experience Replay}
ExpRAG \cite{ferraz2026exprag} stores completed agent trajectories and uses them
twice, as fine-tuning data and as retrieved context during supervised
fine-tuning, so that the model learns to use precedent rather than merely having
precedent available.

\subsection{Attacks on Retrieval}
Corpus poisoning against retrieval-augmented generation is established
\cite{zou2024poisonedrag, zhong2023poisoning}: an adversary who can insert
documents can steer generation. Our source-poisoning condition
(Section~\ref{sec:source}) belongs to this family. The memory-poisoning
condition differs in kind, because the adversary writes to the agent's record of
\emph{its own reasoning}, and the record that eventually does the damage may be
one the agent wrote itself.

\subsection{Fact Verification}
AVeriTeC \cite{schlichtkrull2023averitec} provides real-world claims with
annotated evidence from archived web pages and a four-way verdict including two
uncertainty labels. Those labels matter here: they give a model somewhere to
retreat to, which drives our dominant failure mode (Section~\ref{sec:failures}).

\section{The VeriMem System}

\subsection{Pipeline}
A claim is verified in five steps. First, \textbf{retrieve}: BM25
\cite{robertson2009bm25} over a pooled corpus built from the archived pages
AVeriTeC's annotators cited. Second, \textbf{grade}: a fine-tuned Flan-T5-base
evaluator \cite{chung2022flant5} scores each passage for relevance in
$[-1,1]$, read as $P(\text{yes})$ against $P(\text{no})$ at the first decoded
position, with the claim's grade taken as the maximum. Third, \textbf{route}:
two thresholds, fitted on 150 development claims, select refine, rewrite, or
both. Fourth, \textbf{retrieve experience}: past trajectories, optionally gated.
Fifth, \textbf{verify}: a verdict from either a prompted teacher or a
LoRA-fine-tuned student \cite{hu2021lora, dettmers2023qlora}, after which source
trust is updated and the trajectory is stored.

Four configuration flags control the contributions. With all of them off, the
pipeline must reproduce the baseline exactly. A regression test asserts that
passing the full machinery to a flags-off run produces byte-identical model
calls, and separately that each flag changes behaviour when enabled, so a flag
wired to nothing cannot pass.

\subsection{Evaluator-Gated Experience Retrieval}
A retrieved trajectory is scored by the same evaluator that grades evidence, but
on a different question: does this trajectory's evidence support the verdict it
reached? Our first specification asked instead whether the trajectory was
relevant to the incoming claim, and it was defeated. A poisoned record is
written \emph{about} the target claim, so it is genuinely relevant and the
evaluator scored it highly, correctly. Relevance is not truthfulness. Scoring
evidence against the verdict it is supposed to justify is a question a
fabrication can fail.

\subsection{Bayesian Source Trust}
Each domain carries a $\mathrm{Beta}(\alpha,\beta)$ posterior with
$\text{trust}=\alpha/(\alpha+\beta)$, updated by whether the verdict a
trajectory reached was correct. Trajectories and evidence strips whose domains
fall below a floor of $0.3$ are dropped.

The prior is $\mathrm{Beta}(1,3)$, not $\mathrm{Beta}(1,1)$. A uniform prior
places an unseen domain at $0.5$, above any usable floor, so an adversary
registering fresh domains passes straight through; a model that distrusts only
domains it has already watched fail cannot defend against a new attacker. Under
$\mathrm{Beta}(1,3)$ a stranger starts at $0.25$ and must earn its way in. This
choice turns out to be the whole of the defence's effect
(Section~\ref{sec:trust}).

\section{Experimental Setup}
The official AVeriTeC knowledge store was withdrawn during the project, so we
rebuilt a corpus from the Wayback snapshots in the gold annotations: 764
documents and 11{,}559 passages for development, 3{,}468 and 51{,}471 for
training. Evidence is pooled across claims rather than kept per claim, since a
per-claim pool makes retrieval unable to fail and the baseline meaningless.
Fact-checking domains are blocklisted, and the blocklist resolves archive
wrappers to the underlying publisher; 50 of 961 development URLs were
fact-checkers inside archive URLs and would otherwise have passed.

Teacher and judge are Qwen2.5-7B-Instruct \cite{qwen2024}, served locally. The
student is Qwen2.5-3B-Instruct with QLoRA ($r=16$, $\alpha=32$, 4-bit base). The
evaluator is Flan-T5-base. All runs are single-seed with bootstrap intervals;
see Section~\ref{sec:limitations}. For the attacks, \emph{attack success rate}
is the fraction of targeted claims receiving the adversary's chosen verdict,
reported separately from collateral degradation, since a defence that cuts
attack success by discarding all retrieved experience has not defended anything.

\section{Reproducing CRAG}
\label{sec:crag}

\begin{table}[!t]
\caption{Phase One on AVeriTeC Development, 500 Claims}
\label{tab:crag}
\centering
\begin{tabular}{lcc}
\toprule
System & Macro-F1 & 95\% CI \\
\midrule
No retrieval & 0.324 & [0.288, 0.360] \\
Plain RAG    & 0.354 & [0.319, 0.388] \\
CRAG         & 0.363 & [0.324, 0.404] \\
\midrule
\multicolumn{3}{l}{\emph{Paired differences}}\\
Plain RAG $-$ no retrieval & $+0.029$ & [$-0.016$, $+0.078$] \\
CRAG $-$ plain RAG         & $+0.009$ & [$-0.031$, $+0.051$] \\
CRAG $-$ no retrieval      & $+0.038$ & [$-0.005$, $+0.085$] \\
\bottomrule
\end{tabular}
\end{table}

\subsection{What Reproduced}
The evaluator reaches AUC $0.811$ and accuracy $0.728$ on 751 held-out pairs,
against AUC $0.650$ $[0.611, 0.686]$ and accuracy $0.573$ for an off-the-shelf
cross-encoder: a gain of $+0.161$ AUC from fine-tuning. The reason is specific.
The baseline's accuracy barely clears the majority-class floor of $0.531$ while
its AUC shows real ranking ability, and its mean score on genuinely relevant
passages is $-0.296$. It ranks adequately and decides in the wrong place. Since
CRAG reads two thresholds off this score, a misplaced boundary routes nearly
everything down one branch however well the model ranks; fine-tuning moved the
boundary, not only the ordering. Fitted thresholds beat the best fixed branch by
$0.044$ macro-F1, where guessed thresholds beat it by $0.004$.

\subsection{What Did Not Transfer}
Table~\ref{tab:crag} shows no significant pairwise difference. The null is two
effects cancelling. Refinement keeps the best five \emph{sentences} where plain
RAG passes the best five \emph{passages}, so the verifier sees 1{,}281
characters of evidence instead of 3{,}750, a reduction of 66\%, and it abstains.
CRAG answers ``Not enough evidence'' 145 times against a gold count of 35;
31.0\% of its predictions are an uncertainty label against 14.6\% in gold, while
plain RAG sits at 15.2\%, almost exactly calibrated. Accuracy falls
significantly, $-0.054$ $[-0.096, -0.012]$, while the two rare classes gain, and
macro-F1, weighting all four equally, cancels them.

This is a task-transfer limit rather than a failed reproduction. Refinement was
designed for question answering, where a short answer span suffices and
surrounding context is noise. In fact verification the verdict depends on how
much the evidence \emph{establishes}, and ``Not enough evidence'' is a valid
label to retreat into. Thinning the evidence by two thirds makes that retreat
the easy option.

\section{Reproducing ExpRAG}
\label{sec:exprag}

\begin{table}[!t]
\caption{Four-Row Comparison, 300 Development Claims}
\label{tab:exprag}
\centering
\begin{tabular}{llcc}
\toprule
Adapter & Retrieval & Macro-F1 & Accuracy \\
\midrule
none                & ---  & 0.243 & 0.263 \\
none                & yes  & 0.285 & 0.333 \\
LoRA SFT            & ---  & 0.329 & 0.667 \\
Retrieval-aug.\ SFT & yes  & 0.332 & 0.673 \\
\midrule
\multicolumn{4}{l}{\emph{Filtered targets, for comparison}}\\
LoRA SFT            & ---  & 0.302 & 0.630 \\
Retrieval-aug.\ SFT & yes  & 0.300 & 0.633 \\
\bottomrule
\end{tabular}
\end{table}

Fine-tuning on the agent's own trajectories works: accuracy rises from $0.263$
to $0.667$ and macro-F1 by $+0.086$ (Table~\ref{tab:exprag}).

The central claim does not replicate. ExpRAG argues that training \emph{with}
retrieved trajectories teaches the model to use them, so retrieval-augmented
fine-tuning should beat plain fine-tuning by more than retrieval alone buys.
Decomposed against the prompted baseline, retrieval alone gives $+0.041$,
training alone $+0.086$, and both together $+0.088$, a super-additivity of
$-0.039$. Paired, retrieval-augmented fine-tuning beats plain fine-tuning by
$+0.002$ $[-0.021, +0.026]$: a tight interval, so a measured null rather than an
underpowered test. The prompted-plus-retrieval row is what makes this
falsifiable and is the row most easily omitted; without it, ``retrieval-augmented
fine-tuning is the best row'' cannot be distinguished from retrieval merely
being available at inference.

The mechanism fired. All 1{,}500 training examples carried retrieved context.
The context was empty of signal: measuring each retrieved neighbour's similarity
as a fraction of the claim's own self-score gives a median of $0.203$, with only
21.0\% above $0.30$ and 7.3\% above $0.50$. In roughly four training examples out
of five the two retrieved ``relevant past verifications'' were barely related to
the claim, so the model learned to discount the context. That retrieval still
helps the \emph{prompted} student, $+0.041$, fits the same explanation: it was
never trained to ignore it. ExpRAG's premise requires a store dense enough for
precedents to exist, and 1{,}500 trajectories spanning politics, health and
economics across many countries and years do not contain near-neighbours for
most claims.

One deviation of our own helped. Keeping only the agent's \emph{correct}
verifications starved the minority classes, leaving 5 ``Conflicting evidence''
examples out of 885, because those are exactly the claims the agent gets wrong.
Training toward the claim's gold label instead keeps every labelled trajectory
and raises that count to 106, gaining $+0.027$ $[+0.004, +0.052]$ without
retrieval and $+0.032$ $[+0.009, +0.056]$ with it, both significant.

\section{Memory Poisoning}
\label{sec:c1}

\subsection{Threat Model}
The adversary can write records into the experience store but cannot modify the
model, the retriever, the evidence corpus, or the gold labels. For a target
claim they write trajectories that are close enough to be retrieved beside it
and that confidently assert the verdict they want, citing fabricated evidence.
Every injected record is a syntactically valid trajectory. Fabricated records
paraphrase rather than copy the target claim, since an identical claim is
removed by self-exclusion before it reaches a prompt.

We use two adversaries. The \textbf{naive} one cites domains from AVeriTeC's
published misinformation list. The \textbf{adaptive} one reads the agent's own
history and cites the five domains it trusts most, in our run a major
encyclopedia, a social network, two public-health agencies and a biomedical
archive.

\subsection{Retrieval Is the Vulnerability}
With experience retrieval enabled, attack success climbs monotonically from
$0.220$ at zero poisoning to $0.430$ at twenty percent
(Fig.~\ref{fig:memory}); the intervals at the two ends do not overlap. An agent
with no store sits at $0.160$. Plain retrieval also costs $0.060$ attack success
and $0.040$ macro-F1 at \emph{zero} poisoning: trajectory accuracy is $0.590$,
so clean experience is often wrong experience.

\begin{figure}[!t]
\centering
\includegraphics[width=\columnwidth]{figures/memory_poisoning.png}
\caption{Attack success against poisoned share of the experience store, 100
targets. The no-memory configuration has no store to write to and is shown as a
reference line. Bars are 95\% bootstrap intervals.}
\label{fig:memory}
\end{figure}

\subsection{The Defences Invert}

\begin{table}[!t]
\caption{Attack Success at Twenty Percent Poisoning}
\label{tab:inversion}
\centering
\begin{tabular}{lcc}
\toprule
Configuration & Naive & Adaptive \\
\midrule
Memory (undefended)   & 0.430 & 0.460 \\
$+$ gating            & 0.320 & \textbf{0.310} \\
$+$ trust             & \textbf{0.230} & 0.460 \\
$+$ both              & 0.230 & 0.310 \\
\midrule
No memory             & \multicolumn{2}{c}{0.160} \\
\bottomrule
\end{tabular}
\end{table}

\begin{figure}[!t]
\centering
\includegraphics[width=\columnwidth]{figures/adversary_inversion.png}
\caption{The cells of Table~\ref{tab:inversion} with their intervals. The
reversal between the two defences is wider than the uncertainty on it; the
narrower comparisons between neighbouring bars are not.}
\label{fig:inversion}
\end{figure}

Both defences flatten the curve. But Table~\ref{tab:inversion} shows the result
we consider central: \emph{which} defence does so inverts entirely with the
adversary. Against the naive adversary, trust alone holds attack success at
$0.230$ and keeps zero poisoned records, and gating adds nothing. Against the
adaptive adversary, trust is defeated outright, $0.460$, identical to
undefended, with every poisoned record retained, and gating carries the defence
alone at $0.310$.

Each defence fails against the adversary the other stops. This is a direct
argument that both are necessary; reporting either alone, or only one adversary,
would have supported the opposite conclusion. We also note the limit plainly:
gated memory under the adaptive attack is $0.310$, still nearly double
no-memory's $0.160$. The defences reduce the cost of having a memory. They do
not remove it.

\subsection{What Source Trust Actually Learns}
\label{sec:trust}

\begin{figure}[!t]
\centering
\includegraphics[width=\columnwidth]{figures/trust_divergence.png}
\caption{Trust does not separate known misinformation domains from the rest, and
the ordering it does learn is not a credibility ordering.}
\label{fig:trust}
\end{figure}

The inversion has a single explanation. Fitting trust on all 1{,}500
trajectories gives 954 domains a track record. Mean trust is $0.380$ for the 20
domains on AVeriTeC's misinformation list and $0.375$ for the other 934, a
difference of $0.005$, with the distributions overlapping entirely
(Fig.~\ref{fig:trust}). Among domains with at least ten observations, a social
network is the most trusted at $0.76$, while a leading scientific journal sits at
$0.29$ and a government budget office at $0.23$.

No credibility model produces that ordering. Ours does, because it is not a
credibility model: it credits every domain in a trajectory according to whether
the \emph{verdict} was correct, so what accumulates is the degree to which
citing a domain coincided with being right. Claims evidenced by a budget-office
report are hard; claims evidenced by a social-media post are often easy.

Trust therefore stopped the naive adversary not by recognising misinformation
but because those domains were \emph{unseen}, and $\mathrm{Beta}(1,3)$ places a
stranger below the floor. It is a novelty detector. The adaptive adversary
defeats it by ceasing to be novel, and the two results are one mechanism.

\subsection{Source Poisoning}
\label{sec:source}

\begin{figure}[!t]
\centering
\includegraphics[width=\columnwidth]{figures/source_poisoning.png}
\caption{Poisoning the evidence corpus, and the reason gating cannot defend it:
the evaluator scores fabricated evidence almost as highly as real evidence.}
\label{fig:source}
\end{figure}

Writing fabricated \emph{pages} into the evidence pool is both more effective
and less defensible. Attack success rises from $0.190$ to $0.690$ at five
fabricated pages per target, against memory poisoning's ceiling of $0.430$, and
macro-F1 collapses from $0.374$ to $0.178$. Unlike memory poisoning, this attack
is broadly destructive rather than surgical.

Gating cannot help, and Fig.~\ref{fig:source} measures why rather than asserting
it. A fabricated page is written about the target claim in detail: it \emph{is}
relevant, and the evaluator is right to say so. Mean relevance grade on
fabricated evidence is $0.667$ against $0.740$ on real evidence at one page per
target, narrowing to $0.605$ against $0.643$ at five. Gaps of $0.02$ to $0.07$ on
a $[-1,1]$ scale leave the gate no signal to act on. Trust reduces the attack,
from $0.690$ to $0.440$ at five pages, but against the adaptive adversary it does
nothing at all: $0.620$ against an undefended $0.630$, with every fabricated
strip retained.

\subsection{Propagation}
If a poisoned record makes the agent misjudge claim $A$, and $A$'s verdict is
then stored, does it contaminate a neighbouring claim $B$? The record doing the
harm at that point is not fabricated. It is the pipeline's own output citing real
publishers, so gating sees a coherent trajectory and trust sees ordinary domains.

We paired each of 100 targets with its nearest different claim by similarity
relative to the target's self-score, excluding duplicates, giving 77 pairs. Each
neighbour was verified three times against stores differing \emph{only} in $A$'s
verdict: absent, clean, and corrupted. Of the 28 targets the attack actually
flipped, none propagated. The channel existed, since $A$'s trajectory was
retrieved in 67 of 77 pairs under both the clean and corrupted conditions and in
0 of 77 under absent, so the poisoned verdict reached $B$'s prompt and $B$'s
answer did not move. With 28 trials the honest statement is that one-hop
propagation is below roughly ten percent, not that it is zero. We attribute the
null to the neighbour being recognisably a different claim at a median relative
similarity of $0.23$; direct memory poisoning works because the fabrication
paraphrases its target.

\section{Where the Errors Come From}
\label{sec:failures}
Categorising all 227 errors on the development split by predicates over each
trajectory's own record, over-abstention accounts for 52.9\%: answering ``Not
enough evidence'' where the gold label is definite. False confidence accounts for
20.3\% and polarity flips for 20.3\%.

Only 11 of 227 errors, five percent, had no evidence strip graded above $+0.3$.
In the other 95\% the evaluator judged the evidence relevant and the verifier
still erred, so the failure is in verification rather than retrieval, which
contradicts the intuitive diagnosis. A representative case: a claim graded
$0.9825$ and routed as correct received ``Not enough evidence'', the verifier
reasoning that the evidence said nothing about the specific person named. Both
components were right about their own input. The passages were on topic, which is
what the evaluator scores, and the surviving sentences had lost the specific
fact, which is what the verifier needed. The clearest next change this points at
is passing refined strips \emph{alongside} more of their surrounding passage
rather than instead of it.

\section{Limitations}
\label{sec:limitations}
A human hand-check of 95 sampled relevance pairs puts label noise at $0.289$
when reweighted to the evaluation distribution, with positives the worst class at
$0.360$. A classifier perfect against truth would score about $0.711$ accuracy
against labels that noisy, and our evaluator scores $0.728$, above that ceiling,
which means its errors correlate with the labeller's. Accuracy $0.728$ therefore
reads as agreement with a labelling that is itself 29\% wrong. The cross-encoder
comparison narrows correspondingly: part of the $+0.161$ AUC is a better match to
this labelling rather than to relevance, and the two cannot be separated without
clean labels. The routing thresholds were fitted against the same noisy score.

Runs are single-seed. Every reported number carries a bootstrap interval over the
evaluation set, which answers whether a difference survives a different draw of
claims but does not capture variance from training seeds. A held-out-\emph{topic}
split was planned and not built, so results are on a development subset held out
from the trajectories' source split but not by topic. Retrieval is BM25 rather
than dense embeddings, chosen for determinism when comparing runs that differ
only in a store's contents. Evidence strips were not deduplicated when these runs
were made; repeats cost 5.2\% of evidence slots, measured, and the code was fixed
afterwards without re-running.

\section{Conclusion}
An agent that remembers its own work can be attacked through that memory, and the
attack scales with how much of the store the adversary controls. Two defences
each appear sufficient in isolation and neither is: evaluator gating and source
trust fail against different adversaries, and a study using only one adversary
would have concluded that either alone suffices. The trust model in particular
does not do what it was designed to do. It detects novelty rather than
unreliability, which is both why it stops an attacker using unknown domains and
why it offers nothing against one citing the agent's own favourites. Both
reproductions produced honest negative results, and in each case we located the
mechanism rather than leaving an absence.

\section*{Ethics Statement}
All fabricated pages and trajectories exist only to attack our own system. They
are generated with fixed seeds into a directory excluded from version control,
assembled in memory for a run and discarded; no fabricated content is written
into the evidence corpus on disk, served, or published. Attack domains are drawn
from AVeriTeC's own published list of known misinformation sites, so the study
does not invent new disinformation and does not attribute fabricated claims to a
publisher that has done nothing. We report a genuine weakness in a class of
system now being deployed. We judge that describing the attack surface and two
partial defences is more useful than withholding it, particularly as the stronger
adaptive attack requires only reading a system's own public source history.

\section*{Acknowledgment}
The authors thank the course staff at IIT Bhilai for guidance throughout the
project.

\begin{thebibliography}{10}

\bibitem{yan2024crag}
S.-Q. Yan, J.-C. Gu, Y. Zhu, and Z.-H. Ling, ``Corrective retrieval augmented
generation,'' \emph{arXiv preprint arXiv:2401.15884}, 2024.

\bibitem{ferraz2026exprag}
T. Palmeira Ferraz \emph{et al.}, ``ExpRAG: Retrieval-augmented supervised
fine-tuning over agent experience,'' \emph{arXiv preprint arXiv:2603.18272},
2026.

\bibitem{schlichtkrull2023averitec}
M. Schlichtkrull, Z. Guo, and A. Vlachos, ``AVeriTeC: A dataset for real-world
claim verification with evidence from the web,'' in \emph{Proc. Adv. Neural Inf.
Process. Syst. (NeurIPS) Datasets and Benchmarks Track}, 2023.

\bibitem{zou2024poisonedrag}
W. Zou, R. Geng, B. Wang, and J. Jia, ``PoisonedRAG: Knowledge corruption attacks
to retrieval-augmented generation of large language models,'' in \emph{Proc.
USENIX Security Symp.}, 2024.

\bibitem{zhong2023poisoning}
Z. Zhong, Z. Huang, A. Wettig, and D. Chen, ``Poisoning retrieval corpora by
injecting adversarial passages,'' in \emph{Proc. Conf. Empirical Methods Natural
Language Process. (EMNLP)}, 2023.

\bibitem{hu2021lora}
E. J. Hu \emph{et al.}, ``LoRA: Low-rank adaptation of large language models,''
\emph{arXiv preprint arXiv:2106.09685}, 2021.

\bibitem{dettmers2023qlora}
T. Dettmers, A. Pagnoni, A. Holtzman, and L. Zettlemoyer, ``QLoRA: Efficient
finetuning of quantized LLMs,'' in \emph{Proc. Adv. Neural Inf. Process. Syst.
(NeurIPS)}, 2023.

\bibitem{chung2022flant5}
H. W. Chung \emph{et al.}, ``Scaling instruction-finetuned language models,''
\emph{arXiv preprint arXiv:2210.11416}, 2022.

\bibitem{qwen2024}
Qwen Team, ``Qwen2.5 technical report,'' \emph{arXiv preprint arXiv:2412.15115},
2024.

\bibitem{robertson2009bm25}
S. Robertson and H. Zaragoza, ``The probabilistic relevance framework: BM25 and
beyond,'' \emph{Found. Trends Inf. Retr.}, vol. 3, no. 4, pp. 333--389, 2009.

\end{thebibliography}

\end{document}
